\documentclass[11pt]{article}
\usepackage[a4paper,margin=1.03in]{geometry}
\usepackage{amsmath,amssymb,amsthm,mathtools}
\usepackage[T1]{fontenc}
\usepackage{lmodern,microtype}
\usepackage[colorlinks=true,linkcolor=blue,citecolor=blue,urlcolor=blue]{hyperref}
\usepackage{enumitem}
\setlist{nosep}
\newtheorem{theorem}{Theorem}
\newtheorem{lemma}[theorem]{Lemma}
\newtheorem{proposition}[theorem]{Proposition}
\theoremstyle{remark}\newtheorem{remark}[theorem]{Remark}
\newcommand{\R}{\mathbb R}\newcommand{\N}{\mathbb N}
\newcommand{\C}{\mathbb C}\newcommand{\F}{\mathcal F}
\newcommand{\Om}{\Omega}\newcommand{\dd}{\,\mathrm d}
\newcommand{\vac}{\omega}\newcommand{\fp}{\mathbin{\circledast}}
\DeclareMathOperator{\Ran}{Ran}
\title{A continuous free product system not generated by its tensor core}
\author{Alec Kriebel\\\small Independent researcher\\
\small ORCID: \href{https://orcid.org/0009-0001-9320-500X}{0009-0001-9320-500X}}
\date{Version 1.0, 21 October 2026 UTC\\\small Unrefereed preprint}
\begin{document}
\maketitle
\begin{abstract}
We construct an associative pointed free product system of separable complex
Hilbert spaces whose maximal ordered tensor subsystem does not generate the
whole free system. Its tensor core is the time-ordered finite-word Fock system,
and the closed free subsystem generated by that core is precisely the full
finite-word Fock system. The ambient system has, in addition, nonzero orthogonal
sectors of words of ordinal length below $\omega^2$. A recoverable Gaussian
mixture supplies a sigma-finite measure class invariant under finite-prefix
absorption. Finite monochromatic-run charts turn concatenation into the free
product structure. We prove separability, an explicit continuous Hilbert-field
structure including time zero, and continuity of the full structural maps and
their inverses. This gives a negative answer, in the scalar specialization, to
the universal tensor-core generation question stated by Skeide in 2009. It does
not answer the separate adapted-cocycle or forward tensor-recovery questions.
The argument is analytic; accompanying finite computations do not certify its
infinite-dimensional assertions. This draft is unrefereed and extensively
AI-assisted.
\end{abstract}

\section{The question, conventions, and result}

For pointed complex Hilbert spaces $(G_i,\vac_i)$, write
$G_i^\circ=G_i\ominus\C\vac_i$. Their reduced pointed free product is
\begin{equation}\label{eq:free}
 G_1\fp G_2=\C\vac\ \oplus\!
 \bigoplus_{k\geq1}\ \bigoplus_{\substack{\epsilon\in\{1,2\}^k\\
 \epsilon_j\ne\epsilon_{j+1}}}
 G_{\epsilon_1}^\circ\otimes\cdots\otimes G_{\epsilon_k}^\circ.
\end{equation}
We use the usual associator obtained by flattening centered alternating words.
The ordered tensor inclusion $G_1\otimes G_2\hookrightarrow G_1\fp G_2$
has range
$\C\vac\oplus G_1^\circ\oplus G_2^\circ
 \oplus(G_1^\circ\otimes G_2^\circ)$.

A pointed free product system is a family $(E_t,\vac_t)_{t\geq0}$ with pointed
unitaries
\[
 u_{s,t}:E_s\fp E_t\longrightarrow E_{s+t},\qquad E_0=\C\vac_0,
\]
satisfying associativity and the canonical zero-time identifications. This is
the scalar case of Skeide's Hilbert-bimodule definition
\cite[pp.~529--530]{Skeide2009}; scalar vectors are central and the maps are
coefficient-bilinear. Our convention puts the $s$-factor in the later interval
$(t,s+t)$ and the $t$-factor in $(0,t)$.

For a finite ordered time partition $\pi$ of $t$, let $T_\pi\subset E_t$
be the range of the iterated ordered tensor inclusion. Its tensor core is the
closed subspace intersection $\bigcap_\pi T_\pi$. We compute this intersection
directly and prove it is a tensor subsystem; no general assertion about
intersections of Hilbert-module tensor products is needed. By the free subsystem
generated by a family we mean the smallest closed free subsystem containing it,
equivalently here the closed span of its finite-partition free-product images.

\begin{theorem}\label{thm:main}
There is a pointed free product system of separable complex Hilbert spaces with
the following properties for every $t>0$:
\begin{enumerate}[label=\textup{(\roman*)}]
\item Its maximal ordered tensor subsystem is
\[
 F_t=\C\vac\oplus\bigoplus_{m\geq1}L^2(\Delta_m(t),\dd x),
 \quad\Delta_m(t)=\{t>x_1>\cdots>x_m>0\}.
\]
\item The closed free subsystem generated by $F$ is
\[
 K_t=\C\vac\oplus\bigoplus_{m\geq1}L^2((0,t)^m,\dd x).
\]
\item $K_t\subsetneq E_t$, with a nonzero orthogonal grade-$\omega$ sector.
\end{enumerate}
The system has the explicit continuous-field convention proved in
Section~\ref{sec:continuity}: norm-continuous support-subspace sections, a
continuous vacuum, local triviality at positive time, and continuous full
structural maps and inverses. The same field structure restricts to $F$ and $K$.
\end{theorem}

Skeide sketched the free-product-system viewpoint earlier in the white-noise
discussion of \cite[Section~10.2]{Skeide2006}. The 2009 contribution explicitly
asks whether every free system is generated by its tensor core and gives the
affirmative Fock example \cite[p.~530]{Skeide2009}. Its subsequent construction
from a supplied spatial tensor system and its recovery question run in the
opposite direction. Later references to that forward construction and the
historical limits of this draft are discussed in Section~\ref{sec:scope}.

\section{A recoverable sigma-finite sequence measure}\label{sec:measure}

Let $\lambda$ denote Lebesgue measure on $\R$, let $\varphi$ be the standard
normal density, and let $\gamma$ be the product standard normal probability
measure on $\R^{\N}$, where $\N=\{0,1,\ldots\}$. Put
\[
 \Gamma=\left\{z:2^{-j}z_j\to0,\quad
 \frac1N\sum_{j<N}z_j^2\to1\right\}.
\]
This is Borel and $\gamma(\Gamma)=1$: Gaussian tail bounds and
Borel--Cantelli prove the first condition, and the strong law proves the
second. These are standard probability inputs; see, for example,
\cite{Kallenberg2021}. On $\R\times(0,\infty)\times\Gamma$ use
$\dd a\,\dd\delta/\delta\,\gamma(\dd z)$ and define
\[
 \Theta(a,\delta,z)_j=a+\delta2^{-j}z_j.
\]
Let $C\subset\R^{\N}$ consist of the sequences for which
\begin{equation}\label{eq:recover}
 a(x)=\lim_jx_j\in\R,\qquad
 d(x)^2=\lim_{N\to\infty}\frac1N\sum_{j<N}4^j(x_j-a(x))^2\in(0,\infty).
\end{equation}
All limits and their existence conditions are Borel. The map $\Theta$ is a
Borel bijection onto $C$, with inverse
\[
 a=a(x),\qquad\delta=d(x),\qquad
 z_j=2^j(x_j-a(x))/d(x).
\]
Indeed, these inverse coordinates satisfy exactly the two conditions defining
$\Gamma$. Let $\mu=\Theta_*(\dd a\,\dd\delta/\delta\,\gamma)$.
The sets $|a|\leq N$, $N^{-1}\leq d\leq N$ exhaust $C$ and have measure
$4N\log N<\infty$ for $N\geq2$. Thus $\mu$ is sigma-finite on a standard
Borel space. In particular, this is a mixture over recoverable parameters;
equivalence of distinct fixed-parameter infinite Gaussian laws is never used.

\begin{lemma}[Finite-prefix absorption]\label{lem:prefix}
For each $k\geq0$, prepending $k$ real coordinates is a Borel bijection
$p_k:\R^k\times C\to C$. For every nonnegative Borel $f$,
\begin{equation}\label{eq:prefix-integral}
 \int_C f\dd\mu=
 \int_C\int_{\R^k}f(p_k(v,y))q_k(v;y)\dd v\,\mu(\dd y),
\end{equation}
where $q_0=1$ and
\begin{equation}\label{eq:q}
 q_k(v;y)=\prod_{j=0}^{k-1}
 \frac{\varphi\big((v_j-a(y))/(2^{k-j}d(y))\big)}{2^{k-j}d(y)}.
\end{equation}
These densities are finite and strictly positive, so the product and output
measures are equivalent in both directions. Moreover,
\begin{equation}\label{eq:chain}
 q_{k+\ell}(v,w;y)=q_k(v;p_\ell(w,y))q_\ell(w;y).
\end{equation}
\end{lemma}

\begin{proof}
Deleting $k$ coordinates preserves $a$ and changes $d$ to $2^{-k}d$.
In the second limit in \eqref{eq:recover}, this is the factor $4^{-k}$,
after deletion of finitely many summands and a harmless ratio of Cesaro
lengths. Conversely, a finite prefix preserves the limit and multiplies the
tail scale by $2^k$. These inverse operations are Borel.

In parameter coordinates, the tail has parameters
$(a,2^{-k}\delta,(z_k,z_{k+1},\ldots))$. The measure $\dd\delta/\delta$
is invariant under this scale change, and the Gaussian tail again has law
$\gamma$. Split the first $k$ independent Gaussian coordinates from the tail
and change variables in these $k$ coordinates. Since the original scale is
$2^kd(y)$, the coordinate scales are $2^{k-j}d(y)$ and their densities are
\eqref{eq:q}. Tonelli gives \eqref{eq:prefix-integral}, without requiring
a probability disintegration of a measure of infinite total mass. Each
density integrates to one in $v$ and is everywhere positive and finite.
Finally $d(p_\ell(w,y))=2^\ell d(y)$ makes the first $k$ scales in
\eqref{eq:chain} exactly the first $k$ scales of $q_{k+\ell}$; the remaining
$\ell$ factors agree as well.
\end{proof}

For any Borel bijection $c:(X,\eta)\to(Y,\zeta)$ with two-way measure-class
equivalence, its canonical unitary is
\begin{equation}\label{eq:canonical}
 U_cf=(f\circ c^{-1})\sqrt{\frac{\dd(c_*\eta)}{\dd\zeta}}.
\end{equation}
The Radon--Nikodym chain rule and the positive square root give
$U_{b\circ c}=U_bU_c$. For the prefix map in Lemma~\ref{lem:prefix},
\eqref{eq:prefix-integral} shows that the multiplier from product measure to
$\mu$ is $q_k^{-1/2}$, rather than $q_k^{1/2}$. Although this multiplier
need not be bounded pointwise, the change-of-measure identity makes the map
unitary. This standard $L^2$ transport principle is the only density mechanism
used below.

\begin{lemma}\label{lem:actions}
Common translation preserves $\mu$, and common dilation by $r>0$ satisfies
$\mu(rA)=r\mu(A)$. Their canonical unitaries are strongly continuous on
$L^2(C,\mu)$. For each fixed $c\in\R$, almost every sequence has no
coordinate equal to $c$, has limit different from $c$, and changes sides of
$c$ only finitely many times.
\end{lemma}

\begin{proof}
Translation changes $a$ to $a+h$. Dilation changes $(a,\delta)$ to
$(ra,r\delta)$, contributing the factor $r$ from $\dd a$. In coordinates
$(a,\log\delta,z)$, the unitaries are the usual finite-dimensional translation
and dilation operators tensored with the identity on the Gaussian factor.
Strong continuity follows first on compactly supported continuous functions
of $(a,\log\delta)$ times an $L^2(\gamma)$ vector, then by density. The
events $a=c$ and $x_j=c$ are null by Lebesgue measure and the finite Gaussian
density, respectively. Off their countable union, convergence to $a\ne c$
makes the side eventually constant.
\end{proof}

\begin{lemma}[Rises in an infinite block]\label{lem:rises}
For $\mu$-almost every $x$, both $x_{2j+1}>x_{2j}$ and
$x_{2j+1}<x_{2j}$ occur infinitely often.
\end{lemma}

\begin{proof}
The sign of $x_{2j+1}-x_{2j}$ is that of $z_{2j+1}/2-z_{2j}$.
These variables, for distinct $j$, are independent symmetric nondegenerate
Gaussians. Each sign has probability $1/2$, so Borel--Cantelli gives the claim.
It holds for every $a,\delta$ on the same Gaussian probability-one set, and
Tonelli integrates it against the sigma-finite parameter measure. Restricting
to any measurable interval event preserves this almost-everywhere assertion.
\end{proof}

\section{Ordinal word spaces and fixed-cut factorization}\label{sec:words}

Every $\alpha<\omega^2$ uniquely equals $\omega n+m$ for finite $n,m\geq0$.
Represent a word of that grade by $n$ consecutive infinite sequences and a
final string of $m$ letters. Define
\[
 \Om_\alpha=C^n\times\R^m,\qquad
 \nu_\alpha=\mu^n\times\lambda^m.
\]
For $\alpha=0$, use the singleton empty word of mass one. Let $(\Om,\nu)$
be their countable measured disjoint union and $H=L^2(\Om,\nu)$. These
standard Borel sigma-finite sectors have separable $L^2$ spaces: a countable
generating algebra and a finite-measure exhaustion provide a countable dense
set of simple functions. Hence $H$ is separable. Its vacuum $\vac$ is the
indicator of the empty word. For an open interval $I$, let $\Om_\alpha(I)$
require every letter coordinate to lie in $I$, and put
$E(I)=L^2(\Om(I),\nu)\subset H$. The empty word is always included.

\begin{lemma}[Coordinate-retaining concatenation]\label{lem:concat}
For a fixed finite list $\alpha_1,\ldots,\alpha_r<\omega^2$, concatenation
is a Borel bijection
\[
 c_{\alpha_1,\ldots,\alpha_r}:
 \prod_i\Om_{\alpha_i}\longrightarrow\Om_{\alpha_1+\cdots+\alpha_r}
\]
with two-way measure-class equivalence. The canonical unitaries for these
fixed lists compose associatively under regrouping.
\end{lemma}

\begin{proof}
For two grades, ordinal addition gives
\[
 (\omega n+m)+(\omega p+q)=
 \begin{cases}\omega(n+p)+q,&p>0,\\\omega n+(m+q),&p=0.\end{cases}
\]
If $p>0$, write the two inputs as $(x^1,\ldots,x^n;v)$ and
$(y^1,\ldots,y^p;w)$. The output is
\[
 (x^1,\ldots,x^n,p_m(v,y^1),y^2,\ldots,y^p;w).
\]
All coordinates are retained. The inverse removes the specified first $m$
coordinates of the specified $(n+1)$st block. Lemma~\ref{lem:prefix}
proves bijectivity and measure-class equivalence. If $p=0$, only finite suffix
concatenation occurs. Iteration proves the general statement.

Regrouping retains the same original infinite blocks; every finite string
immediately before one becomes its prefix, and the final finite string remains
the suffix. Thus the coordinate maps agree for any parenthesization. The
chain rule in \eqref{eq:canonical}, explicitly \eqref{eq:chain} for a repeated
prefix absorption, proves equality of the unitaries. This does not assert
cancellation for ordinal addition or a common inverse when the input grades
are allowed to vary.
\end{proof}

\begin{lemma}[Run charts]\label{lem:runs}
For $s,t>0$, concatenation induces a pointed unitary
\[
 E((t,s+t))\fp E((0,t))\longrightarrow E((0,s+t)).
\]
\end{lemma}

\begin{proof}
In each output sector remove all coordinate and block-limit equalities to the
fixed cut $t$. The removed set is null by Lemma~\ref{lem:actions} and
sigma-finiteness of the other factors. Color each remaining coordinate by its
side of $t$. Every infinite block is eventually one color and has finitely many
changes. There are finitely many blocks and a finite suffix, so the entire
ordinal word has finitely many maximal nonempty monochromatic convex runs.
There can be a color change at a limit position $\omega i$, even without a
predecessor there; it contributes at most one additional boundary.

Each run is again an allowed word. If it starts at $\omega i+k$ and reaches
a later block, it consists of the tail from position $k$ in its first block,
any full intervening blocks, and a possible finite prefix in its last block or
finite suffix. The first tail lies in $C$ by Lemma~\ref{lem:prefix}; all full
blocks remain separate. Its grade is below $\omega^2$. A run contained in one
block before its limit is finite. This accounts also for runs crossing several
limit boundaries.

Conversely, concatenating any finite alternating list of nonempty one-color
words yields exactly those maximal runs, after the same null exclusions.
Run count, color pattern, and ordinal run lengths give countably many charts.
For each chart the domain is a Borel product of interval-support sets, and its
image is Borel by the fixed-grade Borel isomorphism in
Lemma~\ref{lem:concat}. Its restrictions preserve measure classes in both
directions. The images of distinct charts are disjoint because maximal runs
are unique, even when their lists have the same ordinal sum. They cover the
output up to a null set. Taking the canonical chart unitaries and their
orthogonal direct sum is exactly \eqref{eq:free}, and is onto.
\end{proof}

\section{The associative free product system}\label{sec:system}

Set $E_t=E((0,t))$ for $t>0$, $E_0=\C\vac$, and $\vac_t=\vac$.
Translate every coordinate in $E_s$ by $t$ and apply
Lemma~\ref{lem:runs} to define $u_{s,t}$. At zero use the canonical
identifications. Translation preserves $\mu$ and finite-letter Lebesgue
measure, so these are pointed unitaries.

\begin{proposition}\label{prop:assoc}
The maps $u_{s,t}$ satisfy the full free-product associativity identity.
\end{proposition}

\begin{proof}
For three positive lengths $r,s,t$, the original factors occupy the intervals
$(s+t,r+s+t)$, $(t,s+t)$ and $(0,t)$. Delete the null sets for the two
fixed cuts. Every word then has a finite maximal-run decomposition in the
three interval colors. The canonical free-product associator identifies both
parenthesizations with the same direct sum of finite words of the three
centered spaces having unequal neighboring colors. Grouping two factors
groups maximal strings avoiding the third; their internal multiplication is
the same concatenation. Lemma~\ref{lem:concat} and the positive
Radon--Nikodym chain rule therefore identify the two operator compositions on
each countable grade chart and hence on the whole space.

The translations also agree: the three original factors are shifted by
$s+t$, $t$, and $0$ in either composition. Common translation commutes with
concatenation, and $q_k(v+h;y+h)=q_k(v;y)$. Thus no extra density factor
arises. Zero parameters and vacua reduce to the canonical identities. The
same reasoning works for every finite regrouping. Only finitely many fixed
cuts are used for any one identity, so a conull set uniform over all real
cuts is neither asserted nor required.
\end{proof}

\section{An explicit continuous-field structure}\label{sec:continuity}

Let $P_t$ be the projection on $H$ given by the indicator of all coordinates
belonging to $(0,t)$ on every nonempty sector and by the identity on the
vacuum; let $P_0$ be the vacuum projection $P_{\rm vac}$. We use the field
$E_t=P_tH$ whose sections are the norm-continuous $H$-valued functions
$x(t)\in P_tH$. This specifies the continuity convention, rather than
leaving an additional bundle axiom implicit.

\begin{proposition}\label{prop:field}
The projections $P_t$ are strongly continuous for $t\geq0$. They define a
separable continuous Hilbert field, locally trivial for $t>0$, with continuous
vacuum section.
\end{proposition}

\begin{proof}
The coordinates of a word are finitely many convergent sequences and a finite
suffix. Their closure is compact; its supremum is a coordinate or one of the
finitely many block limits. For each fixed positive $t$, the event that the
supremum equals $t$ is therefore null, by a countable union of coordinate
equalities and finitely many limit equalities. Hence the interval indicators
converge almost everywhere when $t_j\to t>0$. At zero, every nonempty word
with all coordinates positive has a positive coordinate and is eventually
excluded as $t_j\to0$; a word with a nonpositive coordinate is always excluded.
Dominated convergence against each $|h|^2$ proves strong continuity, without
assuming finite total measure of a sector.

The sections $P_th$ are continuous and pass through every fibre vector. A
countable dense set of $h\in H$ gives a countable fundamental family. These
sections give exactly the stated ambient norm-continuous field by local
approximation at a fibre value.

On grade $\omega n+m$, dilation satisfies
$\nu_\alpha(rA)=r^{n+m}\nu_\alpha(A)$, since each infinite block contributes
one factor from $\dd a$ and each finite letter one from Lebesgue measure.
Its canonical unitary is
\[
 D_rf(x)=r^{-(n+m)/2}f(x/r).
\]
Lemma~\ref{lem:actions} and ordinary finite-dimensional dilation give strong
continuity on each sector. Finite-sector approximation extends it to $H$.
The maps $D_t:E_1\to E_t$ trivialize the field at positive time. No local
triviality at the dimension-one fibre at zero is asserted.
\end{proof}

\begin{proposition}\label{prop:full-cont}
The full free-product structural maps and their inverses are continuous
between the source and target fields, jointly in $(s,t)\in[0,\infty)^2$.
In particular products of continuous ordered tensor sections and elementary
free-word sections are continuous.
\end{proposition}

\begin{proof}
Put $J=H\fp H$ and $A_t=P_t-P_{\rm vac}$ on $H^\circ$.
On the vacuum of $J$ put $Q(s,t)=1$; on each alternating-word summand put
$Q(s,t)$ equal to the tensor product of the appropriate copies of $A_s$ and
$A_t$. Each such tensor product is strongly continuous. Approximation by
finitely many free-word summands proves strong continuity of $Q$, whose range
is the embedded $E_s\fp E_t$. Extend the structure map by zero off that range:
\[
 V(s,t)=u_{s,t}Q(s,t):J\longrightarrow H.
\]
The run-chart construction already proves that these are partial isometries
with $\|V(s,t)\|\leq1$, initial projection $Q(s,t)$, and final projection
$P_{s+t}$.

To prove strong continuity, take first a vector in a fixed alternating pattern
and fixed ordinal input grades, an elementary tensor of fixed ambient vectors.
After $Q(s,t)$, translate the coordinates from the first copy by $t$ and leave
the others fixed. Apply the fixed-list concatenation unitary of
Lemma~\ref{lem:concat}. Strong continuity of $P$ and translation makes the
result jointly norm-continuous. Finite sums of such vectors are dense in
$J$. The uniform norm bound on the already defined $V(s,t)$ passes continuity
to all of $J$. Equivalently, finite pattern-and-grade truncations commute with
$Q$. We do not assume a bounded unrestricted concatenation map across all
patterns with potentially coincident output grades.

This argument includes zero parameters. On a fixed elementary term containing
a centered factor from a fibre tending to zero, $A_th\to0$, so that term
vanishes. A term from the surviving single factor has its usual translated
limit; the vacuum stays fixed. The norm bound handles the full sum.

Strong continuity of $V$ alone would not imply that of its adjoint. Here its
continuous final projections give, for $p=(s,t)\to p_0=(s_0,t_0)$ and fixed
$h\in H$,
\begin{align*}
 \|V(p)^*h-V(p_0)^*h\|^2
 &=\|P_{s+t}h\|^2+\|P_{s_0+t_0}h\|^2\\
 &\quad-2\operatorname{Re}\langle h,V(p)V(p_0)^*h\rangle\longrightarrow0.
\end{align*}
Thus $V^*$ is strongly continuous as well. Restrictions of $V$ and $V^*$ to
the continuous source and target fields prove the assertion. For variable
continuous vectors, the usual norm estimate adds their norm difference to
the fixed-vector estimate. The ordered tensor inclusion is a continuous
restriction of these maps.
\end{proof}

\section{Computing the maximal tensor core}\label{sec:core}

For a finite partition $\pi$ of $(0,t)$, order its cells from latest to
earliest. Its ordered tensor inclusion allows at most one centered word from
each cell, in that order. The run-chart unitary therefore identifies $T_\pi$
with the entire support subspace of words whose interval colors are
nonincreasing in real time along their ordinal positions. Positive finite
density multipliers preserve these support subspaces.

\begin{proposition}\label{prop:intersection}
For every $t>0$, $\bigcap_\pi T_\pi=F_t$ as in Theorem~\ref{thm:main}.
\end{proposition}

\begin{proof}
First intersect just the two-cell partitions at rational $c\in(0,t)$.
Their countably many boundary exceptions have a null union. On its complement,
a word belongs to every such tensor support precisely when its real labels are
nonincreasing throughout its ordinal order. Indeed, if an earlier coordinate
is smaller than a later coordinate, a rational strictly between the two values
detects a low color followed by a high color. Conversely nonincreasing labels
pass every cut. The index set of any one word is countable. Intersections of
these countably many multiplication projections give exactly the $L^2$ support
on that nonincreasing-label set.

Any sector with at least one infinite block has measure zero on this set by
Lemma~\ref{lem:rises}; sigma-finiteness justifies the product null-set
statement. Restricting to interval events cannot revive a null set. In finite
grade $m$, equal-coordinate diagonals are Lebesgue-null, so the remaining
support is $\Delta_m(t)$. Thus the rational-cut intersection is $F_t$, and
the intersection over all partitions is contained in $F_t$.

For the reverse inclusion, fix any one finite partition, with arbitrary real
endpoints. Almost every decreasing finite tuple avoids its boundaries and
has the required cell order. Its full simplex $L^2$ subspace is therefore in
$T_\pi$. Hence $F_t\subset T_\pi$ for each individual $\pi$. This is an
inclusion of closed subspaces and takes no uncountable union of exceptional
null sets. It proves the desired equality.
\end{proof}

\begin{proposition}\label{prop:tensor}
With inherited multiplication, $F$ is a continuous tensor product subsystem
and is maximal among such tensor subsystems.
\end{proposition}

\begin{proof}
Every strictly decreasing finite word in $(0,s+t)$ splits uniquely at $t$
into its later labels followed by its earlier labels, except on the null
boundary. Splitting each simplex according to the number of later labels
identifies its $L^2$ space with the orthogonal sum of the corresponding tensor
products of simplexes. All measures in these finite grades are ordinary
Lebesgue measures, so concatenation has density one. Summing over total word
length proves that $F_s\otimes F_t\to F_{s+t}$ is onto and isometric.
Associativity and continuity follow from the ambient maps.

More explicitly, the fixed global projection on finite decreasing words
commutes with every $P_t$ and with positive dilation. Its fibre restrictions
therefore form the continuous subfield of $F$. Any tensor subsystem with the
inherited maps has its whole time-$t$ fibre in each iterated partition tensor
range $T_\pi$, and hence in their intersection $F_t$. This proves maximality.
\end{proof}

The ordered-simplex realization is the usual time-ordered finite-particle Fock
system. Its tensor-system role is classical; see \cite{Arveson2003,Skeide2006}.
The full finite-word space in the next section is the full Fock space over
$L^2((0,t))$. The extra ordinal sectors, rather than either classical Fock
identification, provide the obstruction to generation.

\section{The generated free hull and the orthogonal obstruction}\label{sec:hull}

Let $K_t$ be the finite-word space in Theorem~\ref{thm:main}. Its fixed global
grade projection commutes with $P_t$ and dilation, so $K$ is a continuous
closed subfield. A finite word has finitely many finite color runs, and a finite
list of finite runs concatenates to a finite word. Both directions of the
fixed-cut unitary therefore restrict to $K_s\fp K_t\simeq K_{s+t}$.
Consequently $K$ is a closed free subsystem and contains $F$.

\begin{proposition}\label{prop:hull}
The closed free subsystem generated by $F$ is exactly $K$, and
$K_t\ne E_t$ for every $t>0$.
\end{proposition}

\begin{proof}
Every finite-partition free word from $F$ has finite total word length, so its
image lies in $K$. Conversely, fix a finite length $m$. Almost every tuple in
$(0,t)^m$ has pairwise distinct coordinates. A partition with rational
interior endpoints can separate their values into distinct cells. For each
assignment of these $m$ cells to the $m$ positions, the full support rectangle
is the image of the alternating singleton sectors of the corresponding
$F$-fibres, under the finite-partition free-product inclusion. Different
positions have different colors, so the word is reduced. Translation gives
the singleton space of each cell. Thus the whole $L^2$ subspace on that
rectangle belongs to the range, not merely its individual coordinate tuples.

There are countably many rational partitions and cell assignments. Their
supports cover all distinct-coordinate tuples. The closed span of their
support ranges is therefore $L^2((0,t)^m)$: a vector orthogonal to all ranges
vanishes on their countable union. Summing over $m$ proves the hull is $K$.
Since $K$ is already closed and orthogonal to every infinite grade, completion
cannot add an infinite-grade vector.

It remains to exhibit a nonzero $L^2$ vector in grade $\omega$ supported in
$(0,t)$. Restrict the parameter region to
\[
 a\in[t/3,2t/3],\qquad \delta\in[t/100,t/50].
\]
It has measure $(t/3)\log2$. Require also
$|z_j|<(t/(3\delta))2^j$ for every $j$. The thresholds are at least
$(50/3)2^j$. Chebyshev's inequality and the countable union bound give a
uniform failure probability at most
\[
 \sum_{j\geq0}\frac{9}{2500\,4^j}=\frac3{625}.
\]
Thus the image $S_t\subset C$ of this Borel region satisfies
\[
 0<\frac{622}{625}\frac t3\log2\ \leq\ \mu(S_t)
 \ \leq\ \frac t3\log2<\infty.
\]
Intersecting with $\Gamma$ changes no probability. The strict inequalities
on the Gaussian coordinates imply $0<a+\delta2^{-j}z_j<t$, so every resulting
word lies in $(0,t)$. The indicator $1_{S_t}$ is a nonzero square-integrable
grade-$\omega$ vector, orthogonal to $K_t$. This proves properness and
completes Theorem~\ref{thm:main}.
\end{proof}

\section{Historical scope and verification limits}\label{sec:scope}

The negative answer here concerns the universal generation assertion in
\cite[p.~530]{Skeide2009}. It is already obtained over $\C$, so it also
precludes an affirmative universal statement over all coefficient algebras.
We make no claim about adapted cocycles, characterization of the associated
noises, or recovery of a supplied spatial tensor system after the forward
tensor-to-free construction. We prove the explicit regularity convention in
Section~\ref{sec:continuity}; no arbitrary additional regularity axiom is
implicitly included.

There is a simple reason why the known forward construction itself cannot
provide this obstruction, even from a non-Fock tensor input. Write
$G(T)=\operatorname{Hull}(T)$ for the free system generated by a supplied
embedded tensor system $T$, and let $C$ be its tensor intersection. For every
partition the tensor multiplication of $T$ identifies $T_t$ with an embedded
ordered product of its smaller fibres; thus $T_t\subset C_t$. Monotonicity
of the closed free hull gives
\[
 G(T)=\operatorname{Hull}(T)\subset
 \operatorname{Hull}(C)\subset G(T).
\]
Consequently the maximal tensor core of $G(T)$ generates all of $G(T)$.
This excludes that particular direct transfer from old tensor examples; it
does not rule out every construction from the earlier literature.

Skeide's later literature cites a work entitled \emph{Free product systems
generated by spatial tensor product systems} as in preparation; see the
bibliography of \cite{SkeideClassification} and
\cite[Section~4 and reference~30]{Skeide2022}. A 2 February 2026 seminar
abstract \cite{Skeide2026} describes the noise-from-free-system and
spatial-tensor-to-generated-free-system constructions. Its full abstract
contains neither the universal generation converse nor the present negative
example. The exact in-preparation manuscript was not recovered, and no
recording, slides or transcript of the talk were inspected. These limits
prevent an inference about unseen manuscript or talk content. In the bounded
primary-source comparison accompanying this paper, no earlier instance of
the particular counterexample was authenticated. This is not a certificate
of global firstness or a claim about worldwide openness in 2026; unpublished,
unindexed or differently formulated work remains outside that conclusion.

Classical full-Fock rectangle approximation, for example
\cite[Section~3 and Theorem~1\textup{(1)}]{AttalNechita2008}, concerns a
specified ambient full Fock space consisting of finite Cartesian-power
sectors. It does not imply density in the extra ordinal sectors of our $E$.
Nonclassical tensor noises and stationary factorizing measure types are
substantial prior constructions \cite{Tsirelson2004,Liebscher2003}.
Their tensor factorization maps and measure-class half-density machinery
should be distinguished from the pointed free-product run maps required
here. A non-Fock tensor remainder alone does not supply these free maps or
the proper free hull of the maximal tensor intersection.

Terminology alone can mislead. Fowler's \emph{Free $E_0$-semigroups}
\cite{Fowler1995} studies the associated Arveson tensor product system; it is
not the arbitrary pointed free-product-system category used here.
Arici and Ge's \emph{Quadratic subproduct systems, free products, and their
$C^*$-algebras} \cite{AriciGe2026} uses a free product operation for discrete
quadratic subproduct systems, a different object and question. Neither
terminological comparison by itself establishes priority. The Gaussian
probability tools, sigma-finite change of measure, positive half-density
transport, and classical Fock constructions are standard inputs rather than
claims of new techniques.

The source comparison record states actual reading scopes rather than
equating keyword screens with complete paper readings. It records 60
overlapping executed query inputs and 26 retrieval requests, of which 24
returned bodies, including navigation pages and an unrelated thesis. These
are search-custody counts, not counts of independent works or whole papers
read. The accessible exact passages support the comparisons just stated.

The theorem is proved by the analytic arguments above. The historical source
packet includes a program checking 4,240 bounded ordinal and Gaussian-scale
identities. Its assertions are removed by optimized Python, so an optimized
run of that historical script is not verification. The accompanying explicitly
guarded controls test finite support combinatorics and density identities; neither
they nor the old program certify infinite Gaussian laws, measure-class
equivalence, strong continuity, or the full theorem. Reproduction documentation
must preserve this distinction.

\paragraph{Authorship, AI assistance, and review status.}
Alec Kriebel is the named human author and independent researcher. AI systems
were used extensively in research planning, proof development and reconstruction,
literature searching and comparison, computational checks, manuscript writing,
and adversarial auditing. These roles were substantive, rather than limited to
copyediting. The accompanying record distinguishes inherited-context AI audits
from any later review without that history; automated agreement is not
independent human peer review. This draft is unrefereed, has received no formal
independent human peer review, and carries no claim of formal proof-assistant
certification. Mathematical and historical claims remain open to correction.

\begin{thebibliography}{99}
\raggedright
\bibitem{Skeide2009}
M. Skeide, Free Product Systems, in \emph{Mini-Workshop: Product Systems and
Independence in Quantum Dynamics}, Oberwolfach Reports \textbf{6} (2009),
528--530. \href{https://doi.org/10.4171/OWR/2009/09}{doi:10.4171/OWR/2009/09}.

\bibitem{Skeide2006}
M. Skeide, \emph{The index of (white) noises and their product systems},
preprint (2006), especially Section~10.2.
\href{https://arxiv.org/abs/math/0601228}{arXiv:math/0601228}.

\bibitem{SkeideClassification}
M. Skeide, \emph{Classification of $E_0$-semigroups by product systems},
preprint (2009), revision of 13 May 2014, bibliography entry [Ske14b].
\href{https://arxiv.org/abs/0901.1798v4}{arXiv:0901.1798v4}.

\bibitem{Skeide2022}
M. Skeide, Spatial Markov semigroups admit Hudson--Parthasarathy dilations,
\emph{SIGMA} \textbf{18} (2022), 071, 8 pp.
\href{https://doi.org/10.3842/SIGMA.2022.071}{doi:10.3842/SIGMA.2022.071}.

\bibitem{Skeide2026}
M. Skeide, \emph{Noises, Spatial Product Systems, and Free Product Systems},
seminar abstract, Saarland University, 2 February 2026.
\href{https://www.uni-saarland.de/fileadmin/upload/lehrstuhl/weber-moritz/ResearchSeminar/Abstract_Michael_Skeide_2602.pdf}{Institutional abstract}.

\bibitem{AttalNechita2008}
S. Attal and I. Nechita, \emph{Discrete approximation of the free Fock space},
preprint (2008).
\href{https://arxiv.org/abs/0810.4070}{arXiv:0810.4070}.

\bibitem{Tsirelson2004}
B. Tsirelson, \emph{Nonclassical stochastic flows and continuous products},
preprint (2004), revision of 18 December 2004.
\href{https://arxiv.org/abs/math/0402431}{arXiv:math/0402431}.

\bibitem{Liebscher2003}
V. Liebscher, \emph{Random sets and invariants for (type II) continuous
tensor product systems of Hilbert spaces}, preprint (2003).
\href{https://arxiv.org/abs/math/0306365}{arXiv:math/0306365}.

\bibitem{Kallenberg2021}
O. Kallenberg, \emph{Foundations of Modern Probability}, third edition,
Springer, 2021. \href{https://doi.org/10.1007/978-3-030-61871-1}{doi:10.1007/978-3-030-61871-1}.

\bibitem{Arveson2003}
W. Arveson, \emph{Noncommutative Dynamics and $E$-Semigroups},
Springer Monographs in Mathematics, Springer, 2003.
\href{https://doi.org/10.1007/978-0-387-21524-2}{doi:10.1007/978-0-387-21524-2}.

\bibitem{Fowler1995}
N. J. Fowler, Free $E_0$-semigroups, \emph{Canadian Journal of Mathematics}
\textbf{47} (1995), 744--785.
\href{https://doi.org/10.4153/CJM-1995-039-0}{doi:10.4153/CJM-1995-039-0}.

\bibitem{AriciGe2026}
F. Arici and Y. Ge, Quadratic subproduct systems, free products, and their
$C^*$-algebras, \emph{Integral Equations and Operator Theory} \textbf{98}
(2026), article 19.
\href{https://doi.org/10.1007/s00020-025-02821-x}{doi:10.1007/s00020-025-02821-x}.
\end{thebibliography}
\end{document}
